% Options for packages loaded elsewhere
\PassOptionsToPackage{unicode}{hyperref}
\PassOptionsToPackage{hyphens}{url}
\documentclass[
]{article}
\usepackage{xcolor}
\usepackage[margin=1in]{geometry}
\usepackage{amsmath,amssymb}
\setcounter{secnumdepth}{-\maxdimen} % remove section numbering
\usepackage{iftex}
\ifPDFTeX
  \usepackage[T1]{fontenc}
  \usepackage[utf8]{inputenc}
  \usepackage{textcomp} % provide euro and other symbols
\else % if luatex or xetex
  \usepackage{unicode-math} % this also loads fontspec
  \defaultfontfeatures{Scale=MatchLowercase}
  \defaultfontfeatures[\rmfamily]{Ligatures=TeX,Scale=1}
\fi
\usepackage{lmodern}
\ifPDFTeX\else
  % xetex/luatex font selection
\fi
% Use upquote if available, for straight quotes in verbatim environments
\IfFileExists{upquote.sty}{\usepackage{upquote}}{}
\IfFileExists{microtype.sty}{% use microtype if available
  \usepackage[]{microtype}
  \UseMicrotypeSet[protrusion]{basicmath} % disable protrusion for tt fonts
}{}
\makeatletter
\@ifundefined{KOMAClassName}{% if non-KOMA class
  \IfFileExists{parskip.sty}{%
    \usepackage{parskip}
  }{% else
    \setlength{\parindent}{0pt}
    \setlength{\parskip}{6pt plus 2pt minus 1pt}}
}{% if KOMA class
  \KOMAoptions{parskip=half}}
\makeatother
\usepackage{color}
\usepackage{fancyvrb}
\newcommand{\VerbBar}{|}
\newcommand{\VERB}{\Verb[commandchars=\\\{\}]}
\DefineVerbatimEnvironment{Highlighting}{Verbatim}{commandchars=\\\{\}}
% Add ',fontsize=\small' for more characters per line
\usepackage{framed}
\definecolor{shadecolor}{RGB}{248,248,248}
\newenvironment{Shaded}{\begin{snugshade}}{\end{snugshade}}
\newcommand{\AlertTok}[1]{\textcolor[rgb]{0.94,0.16,0.16}{#1}}
\newcommand{\AnnotationTok}[1]{\textcolor[rgb]{0.56,0.35,0.01}{\textbf{\textit{#1}}}}
\newcommand{\AttributeTok}[1]{\textcolor[rgb]{0.13,0.29,0.53}{#1}}
\newcommand{\BaseNTok}[1]{\textcolor[rgb]{0.00,0.00,0.81}{#1}}
\newcommand{\BuiltInTok}[1]{#1}
\newcommand{\CharTok}[1]{\textcolor[rgb]{0.31,0.60,0.02}{#1}}
\newcommand{\CommentTok}[1]{\textcolor[rgb]{0.56,0.35,0.01}{\textit{#1}}}
\newcommand{\CommentVarTok}[1]{\textcolor[rgb]{0.56,0.35,0.01}{\textbf{\textit{#1}}}}
\newcommand{\ConstantTok}[1]{\textcolor[rgb]{0.56,0.35,0.01}{#1}}
\newcommand{\ControlFlowTok}[1]{\textcolor[rgb]{0.13,0.29,0.53}{\textbf{#1}}}
\newcommand{\DataTypeTok}[1]{\textcolor[rgb]{0.13,0.29,0.53}{#1}}
\newcommand{\DecValTok}[1]{\textcolor[rgb]{0.00,0.00,0.81}{#1}}
\newcommand{\DocumentationTok}[1]{\textcolor[rgb]{0.56,0.35,0.01}{\textbf{\textit{#1}}}}
\newcommand{\ErrorTok}[1]{\textcolor[rgb]{0.64,0.00,0.00}{\textbf{#1}}}
\newcommand{\ExtensionTok}[1]{#1}
\newcommand{\FloatTok}[1]{\textcolor[rgb]{0.00,0.00,0.81}{#1}}
\newcommand{\FunctionTok}[1]{\textcolor[rgb]{0.13,0.29,0.53}{\textbf{#1}}}
\newcommand{\ImportTok}[1]{#1}
\newcommand{\InformationTok}[1]{\textcolor[rgb]{0.56,0.35,0.01}{\textbf{\textit{#1}}}}
\newcommand{\KeywordTok}[1]{\textcolor[rgb]{0.13,0.29,0.53}{\textbf{#1}}}
\newcommand{\NormalTok}[1]{#1}
\newcommand{\OperatorTok}[1]{\textcolor[rgb]{0.81,0.36,0.00}{\textbf{#1}}}
\newcommand{\OtherTok}[1]{\textcolor[rgb]{0.56,0.35,0.01}{#1}}
\newcommand{\PreprocessorTok}[1]{\textcolor[rgb]{0.56,0.35,0.01}{\textit{#1}}}
\newcommand{\RegionMarkerTok}[1]{#1}
\newcommand{\SpecialCharTok}[1]{\textcolor[rgb]{0.81,0.36,0.00}{\textbf{#1}}}
\newcommand{\SpecialStringTok}[1]{\textcolor[rgb]{0.31,0.60,0.02}{#1}}
\newcommand{\StringTok}[1]{\textcolor[rgb]{0.31,0.60,0.02}{#1}}
\newcommand{\VariableTok}[1]{\textcolor[rgb]{0.00,0.00,0.00}{#1}}
\newcommand{\VerbatimStringTok}[1]{\textcolor[rgb]{0.31,0.60,0.02}{#1}}
\newcommand{\WarningTok}[1]{\textcolor[rgb]{0.56,0.35,0.01}{\textbf{\textit{#1}}}}
\usepackage{graphicx}
\makeatletter
\newsavebox\pandoc@box
\newcommand*\pandocbounded[1]{% scales image to fit in text height/width
  \sbox\pandoc@box{#1}%
  \Gscale@div\@tempa{\textheight}{\dimexpr\ht\pandoc@box+\dp\pandoc@box\relax}%
  \Gscale@div\@tempb{\linewidth}{\wd\pandoc@box}%
  \ifdim\@tempb\p@<\@tempa\p@\let\@tempa\@tempb\fi% select the smaller of both
  \ifdim\@tempa\p@<\p@\scalebox{\@tempa}{\usebox\pandoc@box}%
  \else\usebox{\pandoc@box}%
  \fi%
}
% Set default figure placement to htbp
\def\fps@figure{htbp}
\makeatother
\setlength{\emergencystretch}{3em} % prevent overfull lines
\providecommand{\tightlist}{%
  \setlength{\itemsep}{0pt}\setlength{\parskip}{0pt}}
\usepackage{bookmark}
\IfFileExists{xurl.sty}{\usepackage{xurl}}{} % add URL line breaks if available
\urlstyle{same}
\hypersetup{
  pdftitle={Populations{,} samples{,} and the central limit theorem},
  hidelinks,
  pdfcreator={LaTeX via pandoc}}

\title{Populations, samples, and the central limit theorem}
\author{}
\date{\vspace{-2.5em}}

\begin{document}
\maketitle

\textbf{Inference labs} use the five-step template: Hypothesis →
Visualise → Assumptions → Conduct → Conclude.

\subsection{Learning objectives}\label{learning-objectives}

\begin{itemize}
\tightlist
\item
  Distinguish population parameters from sample estimates and their
  sampling distributions.
\item
  Decompose expected error into bias squared, variance, and irreducible
  noise.
\item
  Demonstrate the central limit theorem by simulation on a strongly
  skewed population.
\end{itemize}

\subsection{Prerequisites}\label{prerequisites}

Weeks 1 and 2.

\subsection{Background}\label{background}

A population is the set of units we wish to learn about. A sample is the
set of units we actually measure. The population has a fixed parameter
--- mean μ, proportion π, whatever --- and the sample has an estimate
--- x̄, p̂ --- which is a random variable because the sample was drawn at
random. Every estimate has a \emph{sampling distribution}: the
distribution of its value across hypothetical repetitions of the
sampling process. Standard errors, confidence intervals, and
\emph{p}-values are all properties of sampling distributions.

Error decomposes into bias and variance. The bias of an estimator is the
systematic offset of its sampling distribution's mean from the
parameter. The variance is the spread of the sampling distribution. Mean
squared error --- bias squared plus variance --- is the natural scalar
summary of how wrong an estimator is on average. A good estimator is not
one with zero bias; it is one with low MSE.

The central limit theorem is the reason the normal distribution is
ubiquitous. No matter how skewed the population (within mild conditions
on variance), the sampling distribution of the mean approaches normal as
\emph{n} grows. The lab makes that convergence visible.

\subsection{Setup}\label{setup}

\begin{Shaded}
\begin{Highlighting}[]
\FunctionTok{library}\NormalTok{(tidyverse)}
\FunctionTok{set.seed}\NormalTok{(}\DecValTok{42}\NormalTok{)}
\FunctionTok{theme\_set}\NormalTok{(}\FunctionTok{theme\_minimal}\NormalTok{(}\AttributeTok{base\_size =} \DecValTok{12}\NormalTok{))}
\end{Highlighting}
\end{Shaded}

\subsection{1. Hypothesis}\label{hypothesis}

Claim: the sampling distribution of the sample mean from an exponential
(strongly right-skewed) population approaches a normal distribution as
\emph{n} increases, with shrinking standard error.

\subsection{2. Visualise}\label{visualise}

Draw from an exponential population and show the distribution of
individual observations.

\begin{Shaded}
\begin{Highlighting}[]
\NormalTok{pop\_mean }\OtherTok{\textless{}{-}} \DecValTok{1} \SpecialCharTok{/}\NormalTok{ rate}
\NormalTok{pop\_sd   }\OtherTok{\textless{}{-}} \DecValTok{1} \SpecialCharTok{/}\NormalTok{ rate}

\NormalTok{pop\_sample }\OtherTok{\textless{}{-}} \FunctionTok{tibble}\NormalTok{(}\AttributeTok{x =} \FunctionTok{rexp}\NormalTok{(}\DecValTok{2000}\NormalTok{, }\AttributeTok{rate =}\NormalTok{ rate))}
\end{Highlighting}
\end{Shaded}

\begin{Shaded}
\begin{Highlighting}[]
\NormalTok{pop\_sample }\SpecialCharTok{|\textgreater{}}
  \FunctionTok{ggplot}\NormalTok{(}\FunctionTok{aes}\NormalTok{(x)) }\SpecialCharTok{+}
  \FunctionTok{geom\_histogram}\NormalTok{(}\AttributeTok{bins =} \DecValTok{40}\NormalTok{, }\AttributeTok{fill =} \StringTok{"grey60"}\NormalTok{, }\AttributeTok{colour =} \StringTok{"white"}\NormalTok{) }\SpecialCharTok{+}
  \FunctionTok{geom\_vline}\NormalTok{(}\AttributeTok{xintercept =}\NormalTok{ pop\_mean, }\AttributeTok{linetype =} \DecValTok{2}\NormalTok{) }\SpecialCharTok{+}
  \FunctionTok{labs}\NormalTok{(}\AttributeTok{x =} \StringTok{"x"}\NormalTok{, }\AttributeTok{y =} \StringTok{"Count"}\NormalTok{)}
\end{Highlighting}
\end{Shaded}

\subsection{3. Assumptions}\label{assumptions}

CLT assumes iid sampling and finite variance. Exponential qualifies. The
speed of convergence is set by how skewed the parent is --- the more
skewed, the larger the \emph{n} needed for normality.

\begin{Shaded}
\begin{Highlighting}[]
\CommentTok{\# population mean, across sample sizes.}
\NormalTok{summarise\_means }\OtherTok{\textless{}{-}} \ControlFlowTok{function}\NormalTok{(n, }\AttributeTok{B =} \DecValTok{2000}\NormalTok{) \{}
\NormalTok{  xbars }\OtherTok{\textless{}{-}} \FunctionTok{replicate}\NormalTok{(B, }\FunctionTok{mean}\NormalTok{(}\FunctionTok{rexp}\NormalTok{(n, }\AttributeTok{rate =}\NormalTok{ rate)))}
  \FunctionTok{tibble}\NormalTok{(}\AttributeTok{n =}\NormalTok{ n,}
         \AttributeTok{bias =} \FunctionTok{mean}\NormalTok{(xbars) }\SpecialCharTok{{-}}\NormalTok{ pop\_mean,}
         \AttributeTok{variance =} \FunctionTok{var}\NormalTok{(xbars),}
         \AttributeTok{mse =} \FunctionTok{mean}\NormalTok{((xbars }\SpecialCharTok{{-}}\NormalTok{ pop\_mean)}\SpecialCharTok{\^{}}\DecValTok{2}\NormalTok{),}
         \AttributeTok{se =} \FunctionTok{sd}\NormalTok{(xbars))}
\NormalTok{\}}
\NormalTok{bv }\OtherTok{\textless{}{-}} \FunctionTok{map\_dfr}\NormalTok{(}\FunctionTok{c}\NormalTok{(}\DecValTok{5}\NormalTok{, }\DecValTok{10}\NormalTok{, }\DecValTok{30}\NormalTok{, }\DecValTok{100}\NormalTok{, }\DecValTok{500}\NormalTok{), summarise\_means)}
\NormalTok{bv}
\end{Highlighting}
\end{Shaded}

\subsection{4. Conduct}\label{conduct}

Plot the sampling distribution of the mean across growing \emph{n}.

\begin{Shaded}
\begin{Highlighting}[]
  \FunctionTok{replicate}\NormalTok{(B, }\FunctionTok{mean}\NormalTok{(}\FunctionTok{rexp}\NormalTok{(n, }\AttributeTok{rate =}\NormalTok{ rate)))}
\ErrorTok{\}}

\NormalTok{means\_df }\OtherTok{\textless{}{-}} \FunctionTok{bind\_rows}\NormalTok{(}
  \FunctionTok{tibble}\NormalTok{(}\AttributeTok{n =} \StringTok{"n = 2"}\NormalTok{,   }\AttributeTok{mean =} \FunctionTok{sim\_means}\NormalTok{(}\DecValTok{2}\NormalTok{)),}
  \FunctionTok{tibble}\NormalTok{(}\AttributeTok{n =} \StringTok{"n = 5"}\NormalTok{,   }\AttributeTok{mean =} \FunctionTok{sim\_means}\NormalTok{(}\DecValTok{5}\NormalTok{)),}
  \FunctionTok{tibble}\NormalTok{(}\AttributeTok{n =} \StringTok{"n = 30"}\NormalTok{,  }\AttributeTok{mean =} \FunctionTok{sim\_means}\NormalTok{(}\DecValTok{30}\NormalTok{)),}
  \FunctionTok{tibble}\NormalTok{(}\AttributeTok{n =} \StringTok{"n = 100"}\NormalTok{, }\AttributeTok{mean =} \FunctionTok{sim\_means}\NormalTok{(}\DecValTok{100}\NormalTok{))}
\NormalTok{) }\SpecialCharTok{|\textgreater{}}
  \FunctionTok{mutate}\NormalTok{(}\AttributeTok{n =} \FunctionTok{factor}\NormalTok{(n, }\AttributeTok{levels =} \FunctionTok{c}\NormalTok{(}\StringTok{"n = 2"}\NormalTok{, }\StringTok{"n = 5"}\NormalTok{, }\StringTok{"n = 30"}\NormalTok{, }\StringTok{"n = 100"}\NormalTok{)))}
\end{Highlighting}
\end{Shaded}

\begin{Shaded}
\begin{Highlighting}[]
\NormalTok{means\_df }\SpecialCharTok{|\textgreater{}}
  \FunctionTok{ggplot}\NormalTok{(}\FunctionTok{aes}\NormalTok{(mean)) }\SpecialCharTok{+}
  \FunctionTok{geom\_histogram}\NormalTok{(}\AttributeTok{bins =} \DecValTok{40}\NormalTok{, }\AttributeTok{fill =} \StringTok{"grey60"}\NormalTok{, }\AttributeTok{colour =} \StringTok{"white"}\NormalTok{) }\SpecialCharTok{+}
  \FunctionTok{facet\_wrap}\NormalTok{(}\SpecialCharTok{\textasciitilde{}}\NormalTok{ n, }\AttributeTok{scales =} \StringTok{"free"}\NormalTok{) }\SpecialCharTok{+}
  \FunctionTok{geom\_vline}\NormalTok{(}\AttributeTok{xintercept =}\NormalTok{ pop\_mean, }\AttributeTok{linetype =} \DecValTok{2}\NormalTok{) }\SpecialCharTok{+}
  \FunctionTok{labs}\NormalTok{(}\AttributeTok{x =} \StringTok{"Sample mean"}\NormalTok{, }\AttributeTok{y =} \StringTok{"Count"}\NormalTok{)}
\end{Highlighting}
\end{Shaded}

The histograms widen (more spread with small \emph{n}) and are strongly
right-skewed at \emph{n} = 2. By \emph{n} = 30 they are close to
symmetric.

\begin{Shaded}
\begin{Highlighting}[]
\NormalTok{bv }\SpecialCharTok{|\textgreater{}} \FunctionTok{mutate}\NormalTok{(}\AttributeTok{expected\_se =}\NormalTok{ pop\_sd }\SpecialCharTok{/} \FunctionTok{sqrt}\NormalTok{(n),}
             \AttributeTok{ratio =}\NormalTok{ se }\SpecialCharTok{/}\NormalTok{ expected\_se)}
\end{Highlighting}
\end{Shaded}

\subsection{5. Concluding statement}\label{concluding-statement}

\begin{quote}
Sampling from an exponential population with mean 1 and SD 1, the sample
mean was unbiased across every \emph{n} simulated
(\textbar bias\textbar{} \textless{} 0.02 in all cases). Its standard
error shrank as 1/√\emph{n} --- ratio of observed to expected SE near 1
throughout --- and its sampling distribution was visibly skewed at
\emph{n} = 2 but approximately normal by \emph{n} = 30. The central
limit theorem is fully operational by modest sample sizes even for a
strongly skewed parent.
\end{quote}

``\emph{n} ≥ 30'' is a rule of thumb, not a theorem. For very skewed
parents, a larger \emph{n} is needed; for symmetric parents, a smaller
\emph{n} may suffice. Always check by simulation when in doubt.

\subsection{Common pitfalls}\label{common-pitfalls}

\begin{itemize}
\tightlist
\item
  Believing the CLT applies to the \emph{data}, not to the \emph{mean of
  the data}. Individual observations stay whatever they were.
\item
  Invoking the CLT to justify normality of a small skewed sample. The
  CLT is about sampling distributions.
\item
  Confusing standard error (a property of the sampling distribution)
  with standard deviation (a property of the data).
\end{itemize}

\subsection{Further reading}\label{further-reading}

\begin{itemize}
\tightlist
\item
  Efron B, Hastie T. \emph{Computer Age Statistical Inference}, chapter
  1.
\item
  Wasserman L. \emph{All of Statistics}, section on limit theorems.
\end{itemize}

\subsection{Session info}\label{session-info}

\begin{Shaded}
\begin{Highlighting}[]

\end{Highlighting}
\end{Shaded}

\subsection{See also --- chapter index}\label{see-also-chapter-index}

\begin{itemize}
\tightlist
\item
  \href{../index.Rmd}{Methods in this chapter}
\item
  \href{../../../ch00-find-your-method.Rmd}{Find your method (Chapter
  0)}
\end{itemize}

\end{document}
